% Generated by officeParser.
\RequirePackage{iftex}
% DVI output (latex, uplatex or platex, then dvipdfmx): a global class option gives every package the dvipdfmx driver.
\def\officeparserdriver{}\ifpdf\else\ifXeTeX\else\def\officeparserdriver{dvipdfmx,}\fi\fi
\documentclass[\officeparserdriver 11pt]{article}
\iftutex
  \usepackage{fontspec}
\else
  \usepackage[T1]{fontenc}
  \usepackage[utf8]{inputenc}
  \usepackage{lmodern}
\fi
\usepackage[paperwidth=595.28pt,paperheight=841.89pt,top=72pt,bottom=72pt,left=72pt,right=72pt]{geometry}
\usepackage{parskip}
\usepackage{graphicx}
\usepackage{array,longtable,multirow}
\usepackage[hyperfootnotes=false]{hyperref}
\hypersetup{colorlinks=true,
  linkcolor=blue,
  urlcolor=blue,
  citecolor=blue,
  pdftitle={DOCX Demo},
  pdfauthor={Kovid Goyal},
  pdfsubject={calibre},
  pdfinfo={Description={Demonstration of DOCX support in calibre}}}
\setcounter{secnumdepth}{-\maxdimen}
\title{DOCX Demo}
\author{Kovid Goyal}
\date{}

\begin{document}

Demonstration of DOCX support in calibre

This document demonstrates the ability of the calibre DOCX Input plugin to convert the various typographic features in a Microsoft Word (2007 and newer) document. Convert this document to a modern ebook format, such as AZW3 for Kindles or EPUB for other ebook readers, to see it in action.

There is support for images, tables, lists, footnotes, endnotes, links, dropcaps and various types of text and paragraph level formatting.

To see the DOCX conversion in action, simply add this file to calibre using the \textbf{“Add Books” }button and then click “\textbf{Convert”. } Set the output format in the top right corner of the conversion dialog to EPUB or AZW3 and click \textbf{“OK”}.

\hfil\break{}

\section{Text Formatting}\label{text-formatting}

\subsubsection{Inline formatting}\label{inline-formatting}

Here, we demonstrate various types of inline text formatting and the use of embedded fonts.\hfil\break{}

Here is some \textbf{bold, }\textit{italic, }bold-italic, underlined and struck out  text. Then, we have a super\textsuperscript{script} and a sub\textsubscript{script}. Now we see some red, green and blue text. Some text with a yellow highlight. Some text in a box. Some text in inverse video.

A paragraph with styled text: subtle emphasis followed by \textbf{strong text }and intense emphasis. This paragraph uses document wide styles for styling rather than inline text properties as demonstrated in the previous paragraph — calibre can handle both with equal ease.

\subsection{Fun with fonts}\label{fun-with-fonts}

This document has embedded the Ubuntu font family. The body text is in the Ubuntu typeface, here is some text in the Ubuntu Mono typeface, notice how every letter has the same width, even i and m. Every embedded font will automatically be embedded in the output ebook during conversion.

\subsection{Paragraph level formatting}\label{paragraph-level-formatting}

You can do crazy things with paragraphs, if the urge strikes you. For instance this paragraph is right aligned and has a right border. It has also been given a light gray background.

For the lovers of poetry amongst you, paragraphs with hanging indents, like this often come in handy. You can use hanging indents to ensure that a line of poetry retains its individual identity as a line even when the screen is too narrow to display it as a single line. Not only does this paragraph have a hanging indent, it is also has an extra top margin, setting it apart from the preceding paragraph.

\section{Tables}\label{tables}

\begin{longtable}{|*{2}{>{\raggedright\arraybackslash}p{\dimexpr(\linewidth-4\tabcolsep-3\arrayrulewidth)/2\relax}|}}
\hline
ITEM & NEEDED \\
\hline
Books & 1 \\
\hline
Pens & 3 \\
\hline
Pencils & 2 \\
\hline
Highlighter & 2 colors \\
\hline
Scissors & 1 pair \\
\hline
\end{longtable}

Tables in Word can vary from the extremely simple to the extremely complex. calibre tries to do its best when converting tables. While you may run into trouble with the occasional table, the vast majority of common cases should be converted very well, as demonstrated in this section. Note that for optimum results, when creating tables in Word, you should set their widths using percentages, rather than absolute units. To the left of this paragraph is a floating two column table with a nice green border and header row.

Now let’s look at a fancier table—one with alternating row colors and partial borders. This table is stretched out to take 100\% of the available width.

{\small\setlength{\tabcolsep}{4pt}
\begin{longtable}{|*{6}{>{\raggedright\arraybackslash}p{\dimexpr(\linewidth-12\tabcolsep-7\arrayrulewidth)/6\relax}|}}
\hline
City or Town & Point A & Point B & Point C & Point D & Point E \\
\hline
Point A & — &  &  &  &  \\
\hline
Point B & 87 & — &  &  &  \\
\hline
Point C & 64 & 56 & — &  &  \\
\hline
Point D & 37 & 32 & 91 & — &  \\
\hline
Point E & 93 & 35 & 54 & 43 & — \\
\hline
\end{longtable}
}

Next, we see a table with special formatting in various locations. Notice how the formatting for the header row and sub header rows is preserved.

\begin{longtable}{|*{4}{>{\raggedright\arraybackslash}p{\dimexpr(\linewidth-8\tabcolsep-5\arrayrulewidth)/4\relax}|}}
\hline
College & New students & Graduating students & Change \\
\hline
 & Undergraduate &  &  \\
\hline
Cedar University & 110 & 103 & +7 \\
\hline
Oak Institute & 202 & 210 & -8 \\
\hline
 & Graduate &  &  \\
\hline
Cedar University & 24 & 20 & +4 \\
\hline
Elm College & 43 & 53 & -10 \\
\hline
Total & 998 & 908 & 90 \\
\hline
\end{longtable}

Source: Fictitious data, for illustration purposes only

Next, we have something a little more complex, a nested table, i.e. a table inside another table. Additionally, the inner table has some of its cells merged. The table is displayed horizontally centered.

\begin{longtable}{|*{2}{>{\raggedright\arraybackslash}p{\dimexpr(\linewidth-4\tabcolsep-3\arrayrulewidth)/2\relax}|}}
\hline
\begin{tabular}{|*{2}{>{\raggedright\arraybackslash}p{\dimexpr(\linewidth-4\tabcolsep-3\arrayrulewidth)/2\relax}|}}
\hline
\multirow{2}{=}{One

Three} & Two \\
\cline{2-2}
 & Four \\
\hline
\end{tabular} & To the left is a table inside a table, with some cells merged. \\
\hline
\end{longtable}

We end with a fancy calendar, note how much of the original formatting is preserved. Note that this table will only display correctly on relatively wide screens. In general, very wide tables or tables whose cells have fixed width requirements don’t fare well in ebooks.

{\scriptsize\setlength{\tabcolsep}{2pt}
\begin{longtable}{|*{14}{>{\raggedright\arraybackslash}p{\dimexpr(\linewidth-28\tabcolsep-15\arrayrulewidth)/14\relax}|}}
\hline
\multicolumn{13}{|>{\raggedright\arraybackslash}p{\dimexpr(\linewidth-28\tabcolsep-15\arrayrulewidth)*13/14+24\tabcolsep+12\arrayrulewidth\relax}|}{December 2007} &  \\
\hline
Sun &  & Mon &  & Tue &  & Wed &  & Thu &  & Fri &  & \multicolumn{2}{>{\raggedright\arraybackslash}p{\dimexpr(\linewidth-28\tabcolsep-15\arrayrulewidth)*2/14+2\tabcolsep+1\arrayrulewidth\relax}|}{Sat} \\
\hline
 &  &  &  &  &  &  &  &  &  &  &  & \multicolumn{2}{>{\raggedright\arraybackslash}p{\dimexpr(\linewidth-28\tabcolsep-15\arrayrulewidth)*2/14+2\tabcolsep+1\arrayrulewidth\relax}|}{1} \\
\hline
 &  &  &  &  &  &  &  &  &  &  &  & \multicolumn{2}{>{\raggedright\arraybackslash}p{\dimexpr(\linewidth-28\tabcolsep-15\arrayrulewidth)*2/14+2\tabcolsep+1\arrayrulewidth\relax}|}{} \\
\hline
2 &  & 3 &  & 4 &  & 5 &  & 6 &  & 7 &  & \multicolumn{2}{>{\raggedright\arraybackslash}p{\dimexpr(\linewidth-28\tabcolsep-15\arrayrulewidth)*2/14+2\tabcolsep+1\arrayrulewidth\relax}|}{8} \\
\hline
 &  &  &  &  &  &  &  &  &  &  &  & \multicolumn{2}{>{\raggedright\arraybackslash}p{\dimexpr(\linewidth-28\tabcolsep-15\arrayrulewidth)*2/14+2\tabcolsep+1\arrayrulewidth\relax}|}{} \\
\hline
9 &  & 10 &  & 11 &  & 12 &  & 13 &  & 14 &  & \multicolumn{2}{>{\raggedright\arraybackslash}p{\dimexpr(\linewidth-28\tabcolsep-15\arrayrulewidth)*2/14+2\tabcolsep+1\arrayrulewidth\relax}|}{15} \\
\hline
 &  &  &  &  &  &  &  &  &  &  &  & \multicolumn{2}{>{\raggedright\arraybackslash}p{\dimexpr(\linewidth-28\tabcolsep-15\arrayrulewidth)*2/14+2\tabcolsep+1\arrayrulewidth\relax}|}{} \\
\hline
16 &  & 17 &  & 18 &  & 19 &  & 20 &  & 21 &  & \multicolumn{2}{>{\raggedright\arraybackslash}p{\dimexpr(\linewidth-28\tabcolsep-15\arrayrulewidth)*2/14+2\tabcolsep+1\arrayrulewidth\relax}|}{22} \\
\hline
 &  &  &  &  &  &  &  &  &  &  &  & \multicolumn{2}{>{\raggedright\arraybackslash}p{\dimexpr(\linewidth-28\tabcolsep-15\arrayrulewidth)*2/14+2\tabcolsep+1\arrayrulewidth\relax}|}{} \\
\hline
23 &  & 24 &  & 25 &  & 26 &  & 27 &  & 28 &  & \multicolumn{2}{>{\raggedright\arraybackslash}p{\dimexpr(\linewidth-28\tabcolsep-15\arrayrulewidth)*2/14+2\tabcolsep+1\arrayrulewidth\relax}|}{29} \\
\hline
 &  &  &  &  &  &  &  &  &  &  &  & \multicolumn{2}{>{\raggedright\arraybackslash}p{\dimexpr(\linewidth-28\tabcolsep-15\arrayrulewidth)*2/14+2\tabcolsep+1\arrayrulewidth\relax}|}{} \\
\hline
30 &  & 31 &  &  &  &  &  &  &  &  &  & \multicolumn{2}{>{\raggedright\arraybackslash}p{\dimexpr(\linewidth-28\tabcolsep-15\arrayrulewidth)*2/14+2\tabcolsep+1\arrayrulewidth\relax}|}{} \\
\hline
\end{longtable}
}

\section{Structural Elements}\label{id_Toc54889875}\label{structural-elements}

Miscellaneous structural elements you can add to your document, like footnotes, endnotes, dropcaps and the like.

\subsection{Footnotes \& Endnotes}\label{id_Toc201580556}\label{footnotes-endnotes}

Footnotes\href{index\_split\_006.html\#note\_1}{\textsuperscript{1}} and endnotes\href{index\_split\_007.html\#note\_2}{\textsuperscript{2}} are automatically recognized and both are converted to endnotes, with backlinks for maximum ease of use in ebook devices.

\subsection{Dropcaps}\label{id_Toc1977424358}\label{dropcaps}

D

rop caps are used to emphasize the leading paragraph at the start of a section. In Word it is possible to specify how many lines of text a drop-cap should use.

\subsection{Links}\label{id_Toc1233048813}\label{links}

Two kinds of links are possible, those that refer to an external website and those that refer to locations inside the document itself. Both are supported by calibre. For example, here is a link pointing to the \href{http://calibre-ebook.com/download}{calibre download page}. Then we have a link that points back to the section on \href{index\_split\_001.html\#id\_Toc2022725662}{paragraph level formatting} in this document.

\subsection{Table of Contents}\label{id_Toc64145348}\label{table-of-contents}

You can see the Table of Contents created by calibre by clicking the Table of Contents button in whatever viewer you are using to view the converted ebook.

\href{index\_split\_000.html\#id\_Toc581531977}{Demonstration of DOCX support in calibre1}

\href{index\_split\_001.html\#id\_Toc2054249818}{Text Formatting1}

\href{index\_split\_001.html\#id\_Toc2137712100}{Inline formatting2}

\href{index\_split\_001.html\#id\_Toc1074133965}{Fun with fonts2}

\href{index\_split\_001.html\#id\_Toc2022725662}{Paragraph level formatting2}

\href{index\_split\_002.html\#id\_Toc28114276}{Tables2}

\hyperref[id_Toc54889875]{Structural Elements4}

\hyperref[id_Toc201580556]{Footnotes \& Endnotes5}

\hyperref[id_Toc1977424358]{Dropcaps5}

\hyperref[id_Toc1233048813]{Links5}

\hyperref[id_Toc64145348]{Table of Contents5}

\href{index\_split\_004.html\#id\_Toc484565143}{Images6}

\href{index\_split\_005.html\#id\_Toc1359965655}{Lists7}

\href{index\_split\_005.html\#id\_Toc1958162433}{Bulleted List8}

\href{index\_split\_005.html\#id\_Toc415190676}{Numbered List8}

\href{index\_split\_005.html\#id\_Toc1093260318}{Multi-level Lists8}

\href{index\_split\_005.html\#id\_Toc1471533984}{Continued Lists8}

\section{Images}\label{images}

Centered images like this are useful for large pictures that should be a focus of attention.

\includegraphics[width={\ifdim 810pt>\linewidth\linewidth\else 810pt\fi},height={\ifdim 810pt>0.9\textheight0.9\textheight\else 810pt\fi},keepaspectratio]{images/image-1.jpg}

There is no analogous technology in ebooks, so the conversion will usually end up placing the image either centered or floating close to the point in the text where it was inserted, not necessarily where it appears on the page in Word.

\section{Lists}\label{lists}

All types of lists are supported by the conversion, with the exception of lists that use fancy bullets, these get converted to regular bullets.

\subsection{Bulleted List}\label{bulleted-list}

\begin{itemize}
\item One
\item Two
\end{itemize}

\subsection{Numbered List}\label{numbered-list}

\begin{enumerate}
\item One, with a very long line to demonstrate that the hanging indent for the list is working correctly
\item Two
\end{enumerate}

\subsection{Multi-level Lists}\label{multi-level-lists}

\begin{enumerate}
\item One
\end{enumerate}

1.1. Two

1.1.1. Three

1.1.2. Four with a very long line to demonstrate that the hanging indent for the list is working correctly.

1.1.3. Five

\begin{enumerate}
\setcounter{enumi}{1}
\item Six
\end{enumerate}

A Multi-level list with bullets:

\begin{itemize}
\item One
\end{itemize}

\begin{itemize}
\item Two
\end{itemize}

\begin{itemize}
\item This bullet uses an image as the bullet item
\end{itemize}

\begin{itemize}
\item Four
\end{itemize}

\begin{itemize}
\item Five
\end{itemize}

\subsection{Continued Lists}\label{continued-lists}

\begin{enumerate}
\item One
\item Two
\end{enumerate}

An interruption in our regularly scheduled listing, for this essential and very relevant public service announcement.

\begin{enumerate}
\setcounter{enumi}{2}
\item We now resume our normal programming
\item Four
\end{enumerate}

\section{Notes}\label{notes}

\begin{description}
\item[{{[}\href{index\_split\_003.html\#back\_note\_1}{←1}{]}}] In paged media, footnotes are usually displayed at the bottom of the text. However, in ebooks, a better paradigm is to make them clickable endnotes that the user can browse at her pleasure. This conversion is handled automatically by calibre.
\end{description}

\begin{description}
\item[{{[}\href{index\_split\_003.html\#back\_note\_2}{←2}{]}}] Endnotes are typically used for longer notes, they remain endnotes when converted into ebook form, except that they have an additional backlink to make it easy to return to the current position after reading the note.
\end{description}

\end{document}
